\ssr{ВВЕДЕНИЕ}


Поиск информации является одной из базовых операций, встречающихся практически в любой прикладной программе: от обработки пользовательских запросов и работы с файлами до построения индексов в базах данных и анализа больших объёмов данных. Производительность поиска напрямую влияет на общее время работы программной системы, поэтому выбор подходящего алгоритма поиска и условий его применения является важной инженерной задачей.

Алгоритмы поиска различаются по требуемым предпосылкам (например, необходимость предварительной сортировки данных), трудоёмкости в среднем и в худшем случаях, затратам памяти, а также устойчивости к особенностям входных данных. В рамках дисциплины «Анализ алгоритмов» изучение алгоритмов поиска позволяет связать теоретические оценки сложности с практическими измерениями времени выполнения и сделать выводы о целесообразности применения того или иного метода в конкретных условиях.

Целью данной лабораторной работы является исследование алгоритмов поиска и сравнение их эффективности на основе теоретических оценок и измерений.

Для достижения поставленной цели необходимо решить следующие задачи:
\begin{enumerate}
	\item рассмотреть основные подходы к поиску элемента по ключу в структурах данных и определить условия применимости соответствующих алгоритмов;
	\item оценить трудоёмкость выбранных алгоритмов (по числу элементарных операций/сравнений) и привести теоретические выводы;
	\item реализовать алгоритмы поиска в виде программы и обеспечить корректность работы на наборе тестов;
	\item провести исследование: измерить время выполнения (и при необходимости дополнительные метрики), проанализировать результаты и сопоставить их с теорией;
	\item сформулировать выводы по результатам сравнения.
\end{enumerate}

